% Generated from paper/full_draft. Edit the Markdown source or migration script.

\section{Introduction}
\label{sec:02-introduction:introduction}

Knowledge-augmented agents improve language-model responses by conditioning generation on external evidence. In many systems this evidence is supplied by a retrieval-augmented generation (RAG) layer, but the underlying control problem is broader: an agent must decide which pieces of structured domain knowledge should enter its limited context. If relevant evidence is not selected, the generator has little chance to recover; if too much evidence is selected, latency, context cost, and distracting noise increase. Practical systems must repeatedly decide how much evidence to retrieve, which route to trust, and when a smaller context is safe.

Dense retrieval is a strong baseline for this problem because it can recover semantically related passages even when the query and evidence do not share the same surface form. However, dense-only retrieval remains a fixed route. It does not explicitly decide when lexical matching is safer, when local cluster structure is informative, when global dense retrieval should remain as a fallback, or when the final retrieved context can be compacted. These questions become more important in vertical-domain knowledge settings, where corpora are often large, terminology-heavy, and shaped by repeated workflows, entities, and user intents.

This paper studies the following hypothesis: vertical-domain knowledge data often exhibits a piecewise relevance structure. Relevant evidence is not uniformly distributed across the embedding space. Instead, it tends to form local regions induced by domain concepts, lexical anchors, task workflows, and user behavior. Pure dense retrieval performs local semantic neighborhood search in one representation space, but it may not fully capture lexical constraints, cluster-level routing, or user-specific relevance. At the same time, geometry alone is not enough: a cluster route that prunes too early can miss the correct evidence, and dense retrieval must remain available as a recall floor.

We propose IntentRoute, a feedback-adaptive route-confidence and context-budget controller. Its central operation is to use confidence to gate a multi-route evidence pool and then apply a calibrated final-context budget. In our retrieval-augmented QA implementation, the route surface includes dense retrieval, BM25 lexical recall, and cluster-local retrieval. A bounded piecewise relevance-manifold hypothesis motivates the local route construction: fixed KMeans/MiniBatchKMeans regions form reproducible arms rather than a claim that geometry alone defines relevance. Trust-weighted LinUCB updates arm values from controlled feedback, turning user feedback into an adaptive estimate of which local routes can be trusted. Low route confidence keeps dense retrieval as a recall floor; the final budget is calibrated separately on the resulting evidence rankings. The name IntentRoute denotes route control conditioned on query intent, local structure, and feedback state; it does not imply a separate intent-classification stage.

Although the control problem appears in memory, graph, tree, and retrieval-backed agents, our empirical validation instantiates it in retrieval-augmented question answering. We therefore use the broader agent framing as motivation and keep the demonstrated claims tied to retrieval-backed evidence selection.

This design separates route construction, confidence estimation, and final context control. Geometry defines a structured local route prior. LinUCB and trust-weighted feedback adapt route confidence over repeated interactions. Dense and BM25 rescue paths protect final recall. A separate calibration stage chooses the final context-budget parameters. A late reranker can improve candidate ordering, while a sentence or prompt compressor can remove redundant text after evidence selection. IntentRoute occupies the upstream route-control and budget-calibration layers and can be composed with those downstream components.

The same distinction separates three cost layers that are often conflated in RAG experiments: the number of source candidates considered during retrieval, the rate at which global dense retrieval is invoked, and the final number of retrieved context tokens sent to the generator. Our main efficiency claim uses the third layer. Earlier candidate-count reductions are useful retrieval-stage diagnostics, but they are not evidence of lower LLM context cost unless the final context itself is reduced. Because these retrieved chunks enter the LLM generator as input tokens, each percentage point of evidence-context reduction translates directly into a proportional per-query inference-cost reduction, a recurring saving that scales with deployment query volume.

We evaluate IntentRoute across seven dataset settings with deliberately tiered evidentiary roles. LoTTE technology/search supplies the full-stack, large-scale quality-efficiency evaluation from 100k to 638k chunks. A separate LoTTE science/search study tests cross-domain transfer at 20k, 100k, and 200k corpus scales. PubMedQA and CovidQA-RAG test biomedical evidence-retrieval transfer under near-ceiling and more discriminative dense baselines, respectively; Banking77 tests feedback adaptation in banking-intent routing; and eManual and CUAD expose duplicate-text, strict chunk-identity, and sparse-ground-truth boundaries in manual and legal retrieval. We do not pool these settings as if their tasks, labels, and evidence strength were interchangeable. Instead, LoTTE technology/search anchors the complete retrieval-and-budget claim, science/search tests cross-domain transfer, and the other datasets test transfer, mechanism, and failure boundaries.

On LoTTE technology/search, we scale from 100k to 638k corpus chunks and compare against dense-only retrieval using \nolinkurl{sentence-transformers/all-MiniLM-L6-v2} with exact cosine search. Our main cost-quality evidence uses a calibration/test context-budget protocol: the final-context policy is selected on calibration queries and frozen before test evaluation. Under this protocol, calibration-eligible operating points at 100k, 200k, and 638k save 6-18\% final evidence-context tokens, while a 400k diagnostic point shows positive frozen-test behavior but does not satisfy the calibration eligibility gate. A normalized five-fold follow-up yields 14.50\% mean saving with no mean Hit change at 400k, but retains substantial policy-selection variance. Across these scales, IntentRoute avoids the larger $\mathrm{Hit@10}$ losses observed under dense-only adaptive truncation, although strict seed-level non-inferiority remains scale-dependent. A conservative confidence-only policy provides a stable baseline, reducing final retrieved context tokens by approximately 4.7-5.3\% across all scales while preserving dense-level query hit. We treat these as bounded operating points rather than universal or statistically significant dominance claims.

We then test whether the mechanism survives stronger alternatives and direct component controls. Matched-backbone BGE-base and E5-base comparisons retain the quality-cost pattern beyond MiniLM, while a BGE quality-first policy shows that the operating frontier can be moved toward higher retrieval quality at a smaller token saving. Random-route, static-geometry, no-feedback, and arm-count ablations separate route confidence from the dense/BM25 rescue surface: geometry and feedback strongly affect route-level reward and cluster hit, but final fused $\mathrm{Hit@10}$ can hide weak routes when rescue remains active. Sentence-MMR and a Selective Context-style prompt-pruning baseline show that compression is a strong shared downstream layer; a cross-encoder reranker improves full-context support but can select longer contexts. A 300-query answer-level evaluation finally compares matched BGE, E5, and SentMMR pipelines under three LLM judges. No judge or shared-key majority finds a statistically significant correctness difference, while faithfulness effects remain method-dependent. LoTTE science/search provides cross-domain evidence; PubMedQA and CovidQA-RAG provide biomedical transfer checks; Banking77 provides an intent-routing feedback check; and eManual and CUAD expose evaluation and data-quality boundaries. Together with feedback-driven hard-case recovery, these settings broaden the transfer, mechanism, and boundary evidence without extending the LoTTE token-saving headline to incomparable tasks.

The contributions of this paper are:

\begin{enumerate}
\item We formulate retrieval-backed evidence selection as a confidence-gated-routing and calibrated-budget problem and introduce a controller that separates route confidence from final evidence-context size while retaining dense retrieval as a recall floor.
\item We operationalize local relevance structure as reproducible cluster arms and trust-weighted feedback as adaptive LinUCB route confidence. Geometry diagnostics, random/static controls, no-feedback ablations, and arm-count sensitivity identify what these components explain and what dense/BM25 rescue masks.
\item We provide frozen calibration/test evidence from 100k to 638k LoTTE technology/search chunks, cross-domain LoTTE science/search replication, matched BGE/E5 backbones, and a tunable BGE quality-first point, separating token reduction from retrieval-quality non-inferiority with paired query-level statistics.
\item We compare against shared sentence and prompt compression plus cross-encoder reranking, showing that IntentRoute is an upstream controller that composes with rather than replaces these downstream layers.
\item We add a 300-query, three-judge answer-level evaluation, controlled feedback recovery, PubMedQA and CovidQA-RAG biomedical transfer checks, Banking77 mechanism checks, and eManual and CUAD boundary analyses to evaluate the controller across seven dataset settings without conflating their evidentiary roles.
\end{enumerate}

The resulting claim is intentionally bounded. IntentRoute is not presented as a universal replacement for dense retrieval, a proof of a relevance manifold, or a statistically superior answer generator. It is a feedback-driven controller that uses local geometry to structure routes, trust-weighted LinUCB to adapt route confidence, and dense retrieval as a recall floor before a separately calibrated final context budget. Reranking and context compression remain compatible downstream layers rather than competing explanations.
